% =====================================================================
%  1bat-ud6-1.tex  —  SESSIÓ 1: Què és la probabilitat? Atzar, certesa i risc
%  (sense preàmbul: s'inclou des de main.tex)
% =====================================================================
\horatitol{Sessió 1 — Què és la probabilitat? Atzar, certesa i risc}%
{Ordenem situacions quotidianes de la impossible a la segura i descobrim que la probabilitat és un número entre 0 i 1 que mesura la incertesa.}

\subsection{Idea clau}

Arrenca l'última unitat del curs, dedicada a una eina que fem servir cada dia sense
adonar-nos: la \textbf{probabilitat}. Un tècnic de serveis a la comunitat ha de valorar
\textbf{riscos} i \textbf{oportunitats} (quina probabilitat té un programa de funcionar,
quin perfil està més exposat a una situació) i, sobretot, ha de llegir críticament les
xifres de risc que circulen pertot. Avui, sense fórmules, fem emergir la idea: la
probabilitat és una \textbf{mesura de la incertesa}.

\begin{idea}
\textbf{La probabilitat com a mesura de la incertesa}
\begin{itemize}[leftmargin=1.4em, itemsep=2pt]
  \item Tota situació d'atzar es pot col·locar en una línia que va de la
        \textbf{impossible} a la \textbf{segura}.
  \item Hi assignem un \textbf{número entre $0$ i $1$}: \quad $0$ = impossible, \quad
        $1$ = segur, \quad $0,5$ = tant sí com no.
  \item En casos senzills i equiprobables, aquest número és
        \[
        \frac{\text{casos favorables}}{\text{casos possibles}}.
        \]
\end{itemize}
\textbf{Tres llenguatges per al mateix número:} una probabilitat es pot dir en
\textbf{fracció} ($\frac{1}{4}$), en \textbf{decimal} ($0,25$) i en \textbf{percentatge}
($25\%$). Són la mateixa cosa.
\end{idea}

\subsection{Exemples}

\begin{exemple}[ordenar de la impossible a la segura]
Col·loquem aquestes situacions a la línia de la probabilitat:
\begin{itemize}[leftmargin=1.4em, itemsep=1pt]
  \item «Que un dau tregui un nombre menor que $7$»: \textbf{segura} ($P=1$): sempre
        passa.
  \item «Que surti cara en llançar una moneda»: \textbf{tant sí com no} ($P=0,5$).
  \item «Que un dau tregui un $6$»: \textbf{poc probable} ($P=\frac{1}{6}\approx 0,17$).
  \item «Que un dau tregui un $7$»: \textbf{impossible} ($P=0$): mai no passa.
\end{itemize}
Ordenades de menys a més probable: el $7$ ($0$), el $6$ ($0,17$), la cara ($0,5$) i el
«menor que $7$» ($1$). Tota la incertesa cap en el segment de $0$ a $1$.
\end{exemple}

\begin{exemple}[casos favorables sobre casos possibles]
En llançar un dau, quina és la probabilitat que surti un nombre \textbf{parell}? Els
casos \textbf{possibles} són $6$ (les cares $1,2,3,4,5,6$); els \textbf{favorables} (que
siguin parells) són $3$ (el $2$, el $4$ i el $6$). Per tant:
\[
P(\text{parell})=\frac{3}{6}=\frac{1}{2}=0,5=50\%.
\]
La mateixa probabilitat, dita de tres maneres: $\frac{1}{2}$, $0,5$ i $50\%$.
\end{exemple}

\subsection{Exercicis resolts}

\begin{exresolt}[situar i quantificar a ull]
Ordena de menys a més probable i assigna a cadascuna un valor aproximat entre $0$ i $1$:
(a) que en treure una carta d'una baralla de $48$ surti d'oros; (b) que demà el sol
surti per l'est; (c) que en tirar un dau surti un nombre més gran que $4$.
\solucio
(b) És \textbf{segura}: $P=1$ (sempre passa).

(c) «Més gran que $4$» són el $5$ i el $6$: $2$ casos de $6$, $P=\frac{2}{6}\approx 0,33$.

(a) En una baralla de $48$ cartes hi ha $12$ d'oros (un pal de cada quatre): $P=\frac{12}{48}=\frac{1}{4}=0,25$.

Ordenades de menys a més probable: (a) $0,25$, (c) $0,33$, (b) $1$.
\resposta{(a) $0,25$; (c) $\approx 0,33$; (b) $1$. Ordre: a, c, b.}
\end{exresolt}

\begin{exresolt}[les tres formes d'una probabilitat]
En una bossa hi ha $20$ boles: $5$ de vermelles i la resta de blanques. Quina és la
probabilitat de treure'n una de vermella? Expressa-la en fracció, decimal i percentatge.
\solucio
Casos favorables: $5$ (les vermelles); casos possibles: $20$ (totes les boles).
\[
P(\text{vermella})=\frac{5}{20}=\frac{1}{4}=0,25=25\%.
\]
\resposta{$\frac{1}{4}=0,25=25\%$.}
\end{exresolt}

\subsection{Exercicis per fer}

\begin{exercicis}
\begin{exlist}
  \item Digues si cada situació és \textbf{impossible} ($P=0$), \textbf{segura} ($P=1$)
        o \textbf{possible} (entre $0$ i $1$):
    \begin{enumerate}[label=\alph*), leftmargin=1.6em, topsep=2pt]
      \item que en tirar un dau surti un $8$;
      \item que en tirar un dau surti un nombre de l'$1$ al $6$;
      \item que en treure una carta surti una figura.
    \end{enumerate}
  \item En tirar un dau, calcula la probabilitat (en fracció, decimal i percentatge) que
        surti: (a) un $3$; (b) un nombre senar; (c) un nombre més gran que $2$.
  \item En una bossa hi ha $10$ boles: $3$ de verdes, $2$ de grogues i $5$ de blaves.
        Calcula la probabilitat de treure'n una de cada color i expressa-les en
        percentatge.
  \item Ordena de menys a més probable, assignant un valor aproximat: «que toqui la
        grossa de Nadal amb un número», «que surti creu en una moneda», «que en tirar un
        dau surti parell».
  \item \textbf{(Interpretació)} Una probabilitat pot valer $1,3$? I $-0,2$? Explica per
        què una probabilitat ha d'estar sempre entre $0$ i $1$.
  \item \textbf{(Raonament)} Per què a un tècnic de serveis a la comunitat li és útil
        pensar en termes de probabilitat? Posa un exemple de l'àmbit social on calgui
        valorar «com de probable» és alguna cosa.
\end{exlist}
\end{exercicis}

% ---------------------------------------------------------------------
\begin{idea}
\textbf{Fitxa de suport} (per tenir a mà durant la unitat):
\begin{itemize}[leftmargin=1.4em, itemsep=1pt]
  \item \textbf{Cas possible:} cadascun dels resultats que poden passar.
  \item \textbf{Cas favorable:} el resultat (o resultats) que ens interessa que passi.
  \item \textbf{Probabilitat} (casos equiprobables) $=\dfrac{\text{favorables}}{\text{possibles}}$.
  \item \textbf{De fracció a decimal:} divideix. \quad \textbf{De decimal a percentatge:}
        multiplica per $100$.
\end{itemize}
\end{idea}

% ---------------------------------------------------------------------
\fullsolucions{Sessió 1}
\begin{solucions}
\begin{sollist}
  \item (a) \textbf{impossible} ($P=0$); (b) \textbf{segura} ($P=1$); (c)
        \textbf{possible} (hi ha figures, però no totes les cartes ho són).
  \item (a) $\frac{1}{6}\approx 0,17=16,7\%$; (b) senars $1,3,5$: $\frac{3}{6}=\frac{1}{2}=0,5=50\%$;
        (c) més grans que $2$ són $3,4,5,6$: $\frac{4}{6}=\frac{2}{3}\approx 0,67=66,7\%$.
  \item Verda $\frac{3}{10}=30\%$; groga $\frac{2}{10}=20\%$; blava $\frac{5}{10}=50\%$
        (sumen $100\%$).
  \item De menys a més probable: «que toqui la grossa» (probabilitat baixíssima,
        gairebé $0$), «creu en una moneda» ($0,5$) i «parell en un dau» ($0,5$). Les
        dues últimes són igual de probables ($0,5$ cadascuna).
  \item No, una probabilitat \textbf{no} pot valer $1,3$ ni $-0,2$. Ha d'estar sempre
        entre $0$ i $1$: $0$ vol dir que el succés és impossible i $1$ que és segur; no
        té sentit «més segur que segur» (per damunt d'$1$) ni «menys que impossible»
        (per sota de $0$).
  \item Resposta oberta. Idea: pensar en probabilitats ajuda a valorar riscos i
        oportunitats amb dades en comptes d'intuïcions. Exemple vàlid: estimar la
        probabilitat que un usuari recaigui sense seguiment, o que un programa de
        prevenció funcioni, per decidir on posar els recursos.
\end{sollist}
\end{solucions}
